\documentclass[11pt]{letter}
\usepackage[margin=1in]{geometry}
\usepackage{hyperref}

\signature{Zijing Wei\\College of Liberal Arts \& Sciences\\University of Illinois at Urbana-Champaign\\zijingw4@illinois.edu}

\address{Zijing Wei\\College of Liberal Arts \& Sciences\\University of Illinois at Urbana-Champaign\\Urbana, IL 61801, USA}

\begin{document}

\begin{letter}{Editor-in-Chief\\IEEE Access}

\opening{Dear Editor,}

We are pleased to submit our manuscript entitled \textbf{``Auditing Engagement Incentives in the Kidfluencer Ecosystem: A Multimodal Weak Supervision Approach''} for consideration as a Regular Paper in \textit{IEEE Access}.

This paper addresses a pressing societal concern at the intersection of artificial intelligence, platform governance, and child welfare. The rapid growth of child influencer (``kidfluencer'') content on YouTube has raised urgent questions about child digital labor exploitation, yet empirical evidence on how platform engagement metrics relate to exploitative content features has remained scarce due to the difficulty of operationalizing and scaling such measurements.

\textbf{Key contributions of this work include:}

\begin{enumerate}
    \item A novel multimodal weak supervision pipeline that combines 33 labeling functions---including GPT-4.1-mini Vision analysis and rule-based heuristics---to detect exploitation risk signals across six literature-grounded dimensions without requiring large-scale manual annotation.
    \item A large-scale observational audit of 4,208 videos across 56 kid-centric YouTube channels, revealing a statistically significant ``engagement premium'' associated with performative labor (+56.0\%), emotional bait (+65.6\%), and privacy violations (+40.3\%).
    \item A multi-annotator validation study ($N=107$ videos) demonstrating strong agreement between our pipeline and trained human annotators (macro-average F1 = 0.911).
    \item Evidence that explicit commercial content does not benefit from this engagement premium ($-3.8\%$, n.s.), suggesting that the platform ecosystem rewards the commodification of children's identity and labor rather than traditional advertising.
\end{enumerate}

We believe this work is well-suited for \textit{IEEE Access} given its interdisciplinary scope spanning AI/NLP methodology, computational social science, and technology policy. The findings have direct implications for platform design, content moderation systems, and ongoing legislative efforts (e.g., the U.S. KIDS Act, Illinois PA 103-0556) aimed at protecting child creators.

We confirm that this manuscript is original, has not been published previously, and is not under consideration for publication elsewhere. All authors have read and approved the final manuscript and agree to its submission to \textit{IEEE Access}.

We look forward to your consideration.

\closing{Sincerely,}

\end{letter}
\end{document}
